\subsection{由数值密度定义的测度}

设 \(g\) 是一个正数值函数，它局部 \(\mu\) -几乎处处（在 \(T\) 中）有定义且局部 \(\mu\) -可积；此时由 \(t\) （使 \(g(t) = +\infty\) 者）组成的集是局部 \(\mu\) -可忽略的，因为 \(g\varphi_{K}\) 是 \(\mu\) -可积的（对每个紧集 \(K\) ）（Ch. IV, §2, No.~3, 命题 7）。现在设 \(g'\) 是一个局部可积的函数，它是正的且有限，并等于 \(g\) 局部 \(\mu\) -几乎处处；置 \(\lambda_t' = g'(t)\varepsilon_t\) 。映射 \(t \mapsto \lambda_t'\) （\(T\) 到 \(\mathcal{M}_+(T)\) 中的）是淡 \(\mu\) -可测的且标量本质可积（或者再说一次，对 \((I,g')\) ——其中 I 是 \(T\) 的恒等映射——是 \(\mu\) -适应的）；积分 \(\nu = \int \lambda_t' d\mu(t)\) 不依赖于特定的函数 \(g'\) （局部几乎处处等于 \(g\) 、用于定义测度 \(\lambda_t'\) 的那一个）。测度 \(\nu\) 由下列条件确定

\[
\int f (t) d \nu (t) = \int f (t) g (t) d \mu (t) \quad \mathrm{for} f \in \mathcal {K} (\mathrm{T}).\tag{1}
\]

现在，如果 \(\theta\) 是一个复测度，且 u 是一个复函数（或取值于 \(\overline{R}\) 中的函数），它局部 \(\theta\) -几乎处处有定义且关于 \(\theta\) 局部可积，那么可以写

\[
\begin{array}{l} {u = g _ {1} - g _ {2} + i (g _ {3} - g _ {4})} \\ {\theta = \mu_ {1} - \mu_ {2} + i (\mu_ {3} - \mu_ {4})} \end{array}\tag{2}
\]

其中 \(\mu_{1} = (\mathcal{R}\theta)^{+}\) ，\(\mu_{2} = (\mathcal{R}\theta)^{-}\) ，\(\mu_{3} = (\mathcal{I}\theta)^{+}\) ，\(\mu_{4} = (\mathcal{I}\theta)^{-}\) （Ch. III, §1, No.~5），而 \(g_{1}, g_{2}, g_{3}, g_{4}\) 有类似的含义；由于 \(|u|\) 是局部 \(|\theta|\) -可积的，诸正函数 \(g_{i}\) （ \(i = 1, 2, 3, 4\) ）中的每一个对每个正测度 \(\mu_{j}\) （ \(j = 1, 2, 3, 4\) ）都是局部可积的，于是映射

\[
f \mapsto \int f (t) u (t) d \theta (t)
\]

在 \(\mathcal{K}(\mathrm{T})\) 上是一个复测度。

定义 2. --- 设 \(\theta\) 是一个复测度，\(u\) 是一个复函数（或取值于 \(\overline{\mathbf{R}}\) 中的函数），它局部 \(\theta\) -几乎处处有定义且局部 \(\theta\) -可积。T 上的复测度 \(f \mapsto \int f u d\theta\) 称为测度 \(\theta\) 与函数 \(u\) 的乘积，或以 \(u\) 为密度（关于 \(\theta\) ）的测度，记作 \(u \cdot \theta\) 。

每个作为一个正测度 \(\mu\) 与一个局部 \(\mu\) -可积函数之乘积的复测度，都称为以 \(\mu\) 为基的测度。

关系式 \(\eta = u \cdot \theta\) 按惯例又写成

\[
d \eta (t) = u (t) d \theta (t).
\]

当 u 处处有定义且连续时，就重新得到 Ch. III, §1, No.~4 中给出的定义。显然，如果 \(u_{1}\) 和 \(u_{2}\) 都是局部 \(\theta\) -可积的，那么 \((u_{1}+u_{2})\cdot\theta=u_{1}\cdot\theta+u_{2}\cdot\theta\) 。类似地，如果 \(\theta_{1}\) 和 \(\theta_{2}\) 是 T 上两个测度，且 u 是一个关于 \(\theta_{1}\) 和 \(\theta_{2}\) 都局部可积的函数，那么 u 关于 \(\theta_{1}+\theta_{2}\) 局部可积，且有 \(u\cdot(\theta_{1}+\theta_{2})=u\cdot\theta_{1}+u\cdot\theta_{2}\) 。

此后我们将限于处处有定义的函数的情形；推广到局部几乎处处有定义的函数，这总是明显的，留给读者。

下面的命题在大多数情形允许把复测度的情形化归为正测度的情形：

命题 2. --- 设 \(\theta\) 是一个复测度，\(u\) 是一个局部 \(\theta\) -可积的复函数；则

\[
\left| u \cdot \theta \right| = \left| u \right| \cdot \left| \theta \right|.\tag{3}
\]

我们先给出一个辅助结果：

引理 1. --- 设 \(\theta\) 是一个复测度，\(f\) 是 \(\overline{\mathcal{L}}_{\mathbf{C}}^{1}(\mathrm{T},\theta)\) 的一个元素。则

\[
\langle | \theta |, | f | \rangle = \sup _ {c \in \mathcal {K} _ {1}} | \langle \theta , c f \rangle | = \sup _ {c \in \mathcal {B} _ {1}} | \langle \theta , c f \rangle |,\tag{4}
\]

其中 \(\mathcal{K}_{1}\) （相应地，\(B_{1}\) ）表示具有紧支集的连续（相应地，Borel）复函数 c 中使 \(|c| \leqslant 1\) 者所成的集。

我们先处理 \(f \in \mathcal{K}(\mathrm{T};\mathbf{C})\) 的情形。显然

\[
\sup _ {c \in \mathcal {K} _ {1}} | \langle \theta , c f \rangle | \leqslant \sup _ {c \in \mathcal {B} _ {1}} | \langle \theta , c f \rangle | \leqslant \langle | \theta |, | f | \rangle
\]

（Ch. IV, §4, No.~2, 命题 2）。另一方面，设 \(g\) 是 \(\mathcal{K}(\mathrm{T};\mathbf{C})\) 中使 \(|g| \leqslant |f|\) 的一个元素；\(g\) 是元素序列 \((g_n)\) （诸元素属于 \(\mathcal{K}(\mathrm{T};\mathbf{C})\) ）的一致极限，这些元素的支集都含在开集 \(\mathrm{U}\) （由 \(t\) 中使 \(f(t) \neq 0\) 者组成）中，并且显然可以假设 \(|g_n| \leqslant |f|\) （对每个 \(n\) ）。置 \(c_n(t) = g_n(t)/f(t)\) （对 \(t \in \mathrm{U}\) ），置 \(c_n(t) = 0\) （对 \(t \notin \mathrm{U}\) ）；则 \(c_n \in \mathcal{K}_1\) ，\(g = \lim_{n \to \infty} c_n f\) ，因此 \(|\langle \theta, g \rangle| = \lim_{n \to \infty} |\langle \theta, c_n f \rangle|\) ，最后

\[
\sup _ {| g | \leqslant | f |, g \in \mathscr {K} (\mathrm{T}; \mathbf {C})} | \langle \theta , g \rangle | \leqslant \sup _ {c \in \mathscr {K} _ {1}} | \langle \theta , c f \rangle |.
\]

注意到这个不等式的第一个端等于 \(\langle|\theta|,|f|\rangle\) （Ch. III, §1, No.~6, 公式 (12)），便得出结论。

其次，用 \(f\) 表示 \(\overline{\mathcal{L}}_{\mathbf{C}}^{1}(\theta)\) 的一个元素，我们来证明 (4) 仍然成立：只须验证这个关系式的三项连续地依赖于 \(f\) （对 \(\overline{\mathcal{L}}_{\mathbf{C}}^{1}(\theta)\) 的拓扑而言），因为它们在稠密子空间 \(\mathcal{K}(\mathrm{T};\mathbf{C})\) 上相合。这由下列不等式立即得出，其中 \(f\) 和 \(f'\) 表示 \(\overline{\mathcal{L}}_{\mathbf{C}}^{1}(\theta)\) 的元素：

\[
| \langle | \theta |, | f | \rangle - \langle | \theta |, | f ^ {\prime} | \rangle | \leqslant \langle | \theta |, | f - f ^ {\prime} | \rangle = \overline {{\mathrm{N}}} _ {1} (f - f ^ {\prime})
\]

\[
| \langle \theta , c f \rangle - \langle \theta , c f ^ {\prime} \rangle | \leqslant \langle | \theta |, | c | | f - f ^ {\prime} | \rangle \leqslant \overline {{{\mathrm{N}}}} _ {1} (f - f ^ {\prime})
\]

对所有 \(c \in B_{1}\) 成立。引理就此证毕。

转到命题 2 的证明，我们把引理应用于函数 uf，这里 h 属于 \(\mathcal{K}_{+}(\mathrm{T})\) 。这就得出：

\[
\langle | \theta |, | u h | \rangle = \sup _ {c \in \mathscr {K} _ {1}} | \langle \theta , c u h \rangle | = \sup _ {c \in \mathscr {K} _ {1}} | \langle u \cdot \theta , c h \rangle | = \langle | u \cdot \theta |, h \rangle .
\]

然而，第一个端也等于

\[
\langle | \theta |, | u | h \rangle = \langle | u | \cdot | \theta |, h \rangle .
\]

因此，两个测度 \(|u| \cdot |\theta|\) 与 \(|u \cdot \theta|\) 相等。

推论. --- 设 \(g_{1}\) 和 \(g_{2}\) 是两个局部 \(\mu\) -可积的数值函数；则

\[
\inf (g _ {1} \cdot \mu , g _ {2} \cdot \mu) = \inf (g _ {1}, g _ {2}) \cdot \mu ; \quad \sup (g _ {1} \cdot \mu , g _ {2} \cdot \mu) = \sup (g _ {1}, g _ {2}) \cdot \mu .
\]

特别地，若 \(g\) 是一个局部 \(\mu\) -可积的数值函数，则

\[
(g \cdot \mu) ^ {+} = g ^ {+} \cdot \mu ; (g \cdot \mu) ^ {-} = g ^ {-} \cdot \mu .
\]

这由命题 2 以及 Ch. II, §1, No.~1 的公式 (6) 立即得出。

